\begin{claim}{C.57}{Proved}{The full local centralizer with exactly two global quotients}
Use C.53's full global family and suppose
$|\mathscr F_O^{\rm all}|=2$ for an open $G$-orbit $O$.
Let $p:X\to O$ be its original universal covering,
$\Delta$ its actual deck group, and
$\mathcal C_O$ the full local centralizer in C.55.
Its group of all source lifts is
\[
 \widetilde{\mathcal C}_O
   =\{q\in\operatorname{Homeo}(X):
                   pq=\bar q p\text{ for some }\bar q\in\mathcal C_O\}.
                                                               \tag{C.57.1}
\]
In the exact pair-coordinate domains of C.45 the
following list is complete.

For the full-plane domain, put
$T_i=\langle b_{i,t}\rangle$ and
$D_i=C_{\operatorname{Homeo}(\R)}(T_i)$.
Each $D_i$ is, after a specified line conjugacy,
the identity group, the integer translation group,
or the real translation group. In those coordinates,
\[
 \widetilde{\mathcal C}_O=D_1\times D_2,\qquad
 \mathcal C_O\cong(D_1\times D_2)/\Delta,\qquad
 O/\mathcal C_O\cong\R^2/(D_1\times D_2).               \tag{C.57.2}
\]
All these are topological-group or quotient-space
identifications through the original covering.
Up to exchange of coordinates, the final quotient is
\[
\begin{array}{c|c|c}
 D_1\times D_2& O/\mathcal C_O &\text{compact}\\ \hline
 1&\R^2&\text{no}\\
 \Z\times1&S^1\times\R&\text{no}\\
 \R\times1&\R&\text{no}\\
 \Z^2&S^1\times S^1&\text{yes}\\
 \R\times\Z&S^1&\text{yes}\\
 \R^2&\{\ast\}&\text{yes}.
\end{array}                                              \tag{C.57.3}
\]

For the half-plane domain, $\Delta=1$ and
$\mathcal C_O=1$, so its quotient is a plane.
For the band domain $D=\{x<y<x+1\}$, put
$e(x,y)=(y,x+1)$. Then
\[
 \widetilde{\mathcal C}_O=\langle e\rangle\cong\Z,\qquad
 \Delta=1\ \text{or}\ \langle e^{2m}\rangle\ (m\geq1),
 \qquad
 \mathcal C_O\cong\langle e\rangle/\Delta.               \tag{C.57.4}
\]
Its quotient $O/\mathcal C_O$ is always the open
M\"obius band. Thus every noncocompact case in this
two-family list has full lifted local centralizer
$1$, $\Z$, or $\R$, with quotient a plane, open
cylinder, line, or open M\"obius band.
This identifies the groups and quotient maps;
it asserts no further contraction at their frontiers.
\end{claim}

\begin{proof}
\emph{The lift group and its one-dimensional factors.}
Every element of (C.57.1) commutes with the based
source path. Indeed its conjugation of that path
is another lift of the same downstairs path,
with the same identity initial value.
The path and its inverse have compact image.
Their generated source group is transitive by
the compact-word, countable-deck Baire argument
of C.49.3; that argument does not require the
generated group to be Lie.
The two based target paths are loops by C.29,
and their groups $T_i$ are transitive on their
real lines, as in C.52.4.

Equation (C.55.5) proves that every local
centralizer element permutes the full two-element
family of global quotient partitions.
Its source lifts consequently have coordinate-pair
or coordinate-exchange form, as proved in C.46.
The actual source deck maps preserve each selected
partition separately: the equation
$q_iF_i=\Phi_i p$ says that $F_i\delta$ and $F_i$
differ by the specified target deck map.
This retains the original deck group inside
(C.57.1).

For any transitive $T\leq\operatorname{Homeo}^+(\R)$,
its full centralizer $D$ acts freely. An element
fixing one point fixes its entire transitive
$T$-orbit. A decreasing real-line homeomorphism
has a fixed point, so no such element belongs
to $D$. The line version of C.30, used with
the closed transitive group $\overline T$,
makes $D$ locally compact, free and proper
in its actual topology.
The C.48 argument makes it a Lie group
of dimension at most one.

In dimension one its identity component
has open orbits, hence is transitive on
the connected line. Freeness excludes
additional components. Evaluation gives
the standard real translation action.
In dimension zero the orbit of a point
is a closed discrete ordered set.
A nonidentity free increasing action
has points of that orbit on both sides.
The element carrying the point to its
nearest successor generates the group:
its positive iterates cannot have a
finite accumulation point, by properness,
and there is no orbit point between
consecutive iterates. Freeness identifies
every group element with the unique
iterate having its value at that point.
An interpolation on one fundamental
interval and its translates conjugates
the action to integer translations.
The remaining case is the identity group.
This proves the asserted exact alternatives
for every $D_i$, including their topology.

\emph{The plane: the graph obstruction and every descending pair.}
Use the C.45 coordinate map $X\to\R^2$.
A coordinate exchange has the form
$q(x,y)=(a(y),b(x))$. Commutation with every
source generator would imply
\[
                   a b_{2,t}=b_{1,t}a
                                      \quad(0\leq t\leq1).
                                                               \tag{C.57.5}
\]
The proper nonempty graph
$\{(a(y),y):y\in\R\}$ is then preserved
by every generator, hence by the
transitive generated source group.
This is impossible. The argument allows
either orientation of $a$; no exchange
is omitted.

Every lift in (C.57.1) is therefore a
coordinate pair $(a(x),b(y))$.
Commutation with the path is exactly
$a\in D_1$, $b\in D_2$.
Thus the lift group is contained in
$P=D_1\times D_2$. In particular
$\Delta\subset P$.
The preceding one-dimensional
classification makes $P$ abelian.
Conversely every pair in $P$ commutes
with the original lifted path and
with $\Delta$. It therefore descends
through the original $p$ to a
homeomorphism of $O$ commuting with
every $f_t|_O$. Hence it belongs to
(C.57.1). This proves equality with $P$.

The source covering gives
$O\cong\R^2/\Delta$ and the exact
surjection $P\to\mathcal C_O$ has kernel
$\Delta$. Its quotient topology is
the actual compact-open topology:
the covering-chart comparison for
maps and inverses in C.48 proves
this for the group of all lifts,
and the group here is exactly that
group by (C.57.1).
The two maps from $X$ to the orbit
quotient have exactly the same fibers:
\[
      X\longrightarrow O\longrightarrow O/\mathcal C_O,
                   \qquad X\longrightarrow X/P.        \tag{C.57.6}
\]
Both are open quotient maps. Their
induced bijection is consequently a
homeomorphism, proving (C.57.2).
For completeness, the deck group
in these formulas is $1$ or $\Z$:
C.44 already proves that the source
of an actual nontrivial global quotient
is a plane or an open annulus.
Its particular embedding in $P$
remains the original deck embedding,
and is not discarded in the formula
for $\mathcal C_O$.

Choose the line conjugacies just
constructed for $D_1,D_2$ independently.
They give the six rows of (C.57.3).
The quotient maps, in the order
displayed there, are
$(x,y)$, $([x],y)$, $y$,
$([x],[y])$, $[y]$, and the constant
map. Here brackets mean the quotient
by the specified integer translations.
These formulas explicitly identify
the quotient spaces and their
compactness, independently of which
allowed source deck subgroup occurs.

\emph{The half-plane: no pair and no exchange survives.}
The normalized source path is
$(g_t(x),g_t(y))$ on $x<y$.
Write $T=\langle g_t\rangle$.
It acts transitively on ordered pairs
in this half-plane.
For any fixed $x$, its stabilizer
$T_x$ is transitive on each of
the two open rays about $x$:
use ordered pairs $(x,y)$ on the
right and $(y,x)$ on the left.
Consequently
$\operatorname{Fix}(T_x)=\{x\}$.

The coordinate functions of any
pair or exchange commuting with
this diagonal path centralize $T$.
They must be increasing by the
free-centralizer argument above.
In particular a decreasing exchanged
pair, which could otherwise preserve
the ordered half-plane, is excluded
by commutation and transitivity.
An increasing exchanged pair
$(a(y),b(x))$ cannot preserve the
half-plane: for fixed $x$, taking
$y\to+\infty$ makes $a(y)>b(x)$.

For a nonexchanged increasing pair
$(a(x),b(y))$, bijection of the
vertical ray over $x$ onto the
vertical ray over $a(x)$ gives
$b(x)=a(x)$.
Thus it is $(c,c)$ with $c$ centralizing
$T$. The point $c(x)$ is fixed by
$T_x$, so $c(x)=x$ for every $x$.
The whole lift group is therefore
trivial. It contains $\Delta$,
which is also trivial.
The source is the half-plane itself,
homeomorphic to a plane by
$(x,y)\mapsto(x,\log(y-x))$, and
its full local centralizer is trivial.

\emph{The band: the primitive period and the exact quotient.}
Now the source path is diagonal
on $x<y<x+1$, and $T$ commutes with
the unit translation $\tau(x)=x+1$.
Transitivity on the band makes
$T_x$ transitive on $(x,x+1)$.
Commutation with $\tau$ gives
transitivity on each
$(x+n,x+n+1)$, while every $x+n$
is fixed. Hence
\[
                         \operatorname{Fix}(T_x)=x+\Z.
                                                               \tag{C.57.7}
\]
For $c$ in the full line centralizer
of $T$, the point $c(x)$ belongs
to (C.57.7), so $c(x)=x+n$.
Both $c$ and $\tau^n$ commute with
the transitive group $T$.
Agreement at one point makes them
equal everywhere. Thus the full
line centralizer is exactly
$\langle\tau\rangle$.
In particular the given unit period
is primitive; a smaller positive
translation would contradict the
transitivity of $T_x$ on $(x,x+1)$.

The two coordinate functions of a
lifted local centralizer element
are therefore integer translations.
A nonexchanged pair
$(x+m,y+n)$ preserves the band
exactly when $m=n$, giving $e^{2m}$.
An exchanged pair
$(y+m,x+n)$ preserves it exactly
when $n=m+1$, giving $e^{2m+1}$.
Decreasing coordinate functions
are excluded by the same
free-centralizer argument.
This proves
$\widetilde{\mathcal C}_O\subset\langle e\rangle$.

Every source deck map preserves
the plane orientation induced
from the orientable open surface
$O$. The map $e$ reverses orientation,
whereas $e^2(x,y)=(x+1,y+1)$.
Thus $\Delta\subset\langle e^2\rangle$,
and it has one of the forms stated
in (C.57.4).
Every power of $e$ commutes with
the diagonal path and with this
deck group, so it descends.
The reverse inclusion in (C.57.4)
and the actual topological quotient
follow exactly as in the plane case.

Finally the explicit coordinates
\[
          (u,v)=(x+y,\ y-x-\tfrac12)
        \in\R\times(-\tfrac12,\tfrac12)                 \tag{C.57.8}
\]
send $e$ to $(u,v)\mapsto(u+1,-v)$.
Their inverse is
$x=(u-v-\tfrac12)/2$,
$y=(u+v+\tfrac12)/2$.
The quotient by $\langle e\rangle$
is therefore the open M\"obius
band, through this exact
glide-reflection quotient map.
For $\Delta=\langle e^{2m}\rangle$,
the source quotient identifies
$u$ modulo $2m$ and is the open
annulus; its local centralizer
is the cyclic group of order
$2m$. With $\Delta=1$ the source
is a plane and that centralizer
is infinite cyclic.
In both cases the final operator
quotient is the same open
M\"obius band.

The table and the two other
domains now give the last
classification in the claim.
All these coordinate comparisons
identify existing groups, covers,
and quotients. They do not
identify an arbitrary intrinsic
isotopy with an ambient extension
at their noncompact ends.
\end{proof}
